\documentclass[10pt,a4paper]{amsart}
\usepackage[margin=19mm]{geometry}
\usepackage{amsmath,amssymb,mathtools}
\usepackage[scr=rsfs,cal=euler]{mathalfa}
\usepackage{xfrac}
\usepackage[unicode,hidelinks,bookmarks=false]{hyperref}
\newcommand{\ie}{i.e.\ }
\newcommand{\cf}{cf.\ }
\newcommand{\resp}{respectively\ }
\newcommand{\Subsection}{\subsection}
\providecommand{\nobreakdash}{\nobreak\hbox{-}}
\setlength{\emergencystretch}{2em}
\multlinegap0pt
\usepackage{bm}
\usepackage{amssymb}
\usepackage{float}
\usepackage{cancel}
\usepackage[normalem]{ulem}
\newtheorem{theorem}{Theorem}
\newcommand{\beq}{\begin{eqnarray}}
\newcommand{\eeq}{\end{eqnarray}}
\title{Commutant lifting and interpolation on quotients of bounded symmetric domains}
\author{Milan Kumar Mal}
\date{Source: arXiv:2606.06051v1 (CC BY 4.0)}
\begin{document}
\maketitle

\noindent\emph{Source and licence.} Commutant lifting and interpolation on quotients of bounded symmetric domains. Version arXiv:2606.06051v1 (CC BY 4.0), distributed by arXiv under the Creative Commons Attribution 4.0 International licence, http://creativecommons.org/licenses/by/4.0/, as recorded in the arXiv licence metadata for this exact version and shown on the abstract page https://arxiv.org/abs/2606.06051v1. Full licence notice: \texttt{licenses/arxiv-2606.06051-CC-BY-4.0.txt}.\par\smallskip

\noindent\textit{Editorial scope.} Counted unit: the article's theorem 1st (source0.tex lines 841--853). The article's complete original proof of it is delivered with it (source0.tex lines 856--879, one \texttt{proof} environment of 24 lines). No mathematics is written, repaired or re-derived: the statement and the proof are the article's own lines, in the article's own notation, and no retained line is substituted or re-flowed.  Retained uncounted as the source-local context the proof needs: lines 582--584.  Nothing was omitted from any retained span, so the package carries no omission record.  No retained line contains a citation, so the wrapper carries no bibliography and no bibliography entry.  No retained line contains a figure environment, an image-inclusion command or drawing code, so no image is delivered and the package declares no asset dependency.  Editorial formatting only, in addition: an AMS \texttt{amsart} preamble that restates the article's own declarations and macros used by the retained lines, namely the environments \texttt{theorem} on separate counters, together with the standard packages those lines need.  The retained environments print in this document as Theorem 1 in order of appearance, whereas the article prints the same items with the numbering of its own full document.  The complete native source of the article as distributed in the arXiv e-print for this exact version is bundled unchanged as \texttt{sources/202153/source0.tex}.
\begin{theorem}\label{1st characterization}
    Let $\mathcal{Q}\subseteq  H^2(\partial\boldsymbol\theta(\Omega))$ be a quotient module and let $X\in \mathcal{B}_1(\mathcal{Q})$ be a contractive module map. Then $X$ admits a Schur-class lift if and only if the functional $X_{\mathcal{Q}}: \left(\mathcal{M}_{\mathcal{Q}}, \|\cdot\|_{L^1({\partial\boldsymbol\theta(\Omega)})}\right) \to \mathbb C$ defined by \beq \label{eq 11} X_{\mathcal{Q}}(f)=\int_{\partial\boldsymbol\theta(\Omega)}\psi f d\mu_{\boldsymbol\theta},\qquad f\in \mathcal{M}_{\mathcal{Q}}\eeq is contractive, where $\psi=XP_{\mathcal{Q}}1.$
\end{theorem}

\begin{theorem}\label{interpolation 1st}
    Let $\mathcal{Z}=\{\boldsymbol \zeta_1,\ldots,\boldsymbol \zeta_p\}\subseteq \boldsymbol \theta(\Omega)$ be a set of $p$ distinct points, and let $w_1,\ldots,w_p\in \mathbb D.$
Define the scalars $c_1,\ldots,c_p$ via the Gram matrix inverse:
    $$\begin{bmatrix}
            c_1 \\  \vdots\\ c_p 
        \end{bmatrix}= {\begin{bmatrix}
            \mathbb S_{\boldsymbol \theta(\Omega)}(\boldsymbol \zeta_i, \boldsymbol \zeta_j)
        \end{bmatrix}}_{p\times p}^{-1} \begin{bmatrix}
            w_1 \\ \vdots \\w_p
        \end{bmatrix}.$$
    Set $\mathcal{M}_{\mathcal{Q}_{\mathcal{Z}}}=\mathcal{Q}_{\mathcal{Z}}^{conj} \dotplus \mathcal{M}_{\boldsymbol \theta(\Omega)}\dotplus H^2_{0}(\partial\boldsymbol \theta(\Omega))$ and $\psi_{\mathcal{Z}, \mathcal{W}}=\sum_{i=1}^p c_i \mathbb S_{\boldsymbol \theta(\Omega)}(\cdot, \boldsymbol \zeta_i). $ 
    Then there exists $\varphi\in \mathcal{S}(\boldsymbol \theta(\Omega))$ satisfying $\varphi(\boldsymbol \zeta_i)=w_i$ for all $i=1,\ldots,p$ if and only if the functional $X_{\mathcal{Q_{\mathcal{Z}}}}: \left(\mathcal{M}_{\mathcal{Q}_{\mathcal{Z}}}, \|\cdot \|_{L^1(\partial\boldsymbol \theta(\Omega))}\right) \to \mathbb C$ defined by $$X_{\mathcal{Q_{\mathcal{Z}}}} f=\int_{\partial\boldsymbol \theta(\Omega)} \psi_{\mathcal{Z}, \mathcal{W}} fd\mu_{\boldsymbol \theta}, \qquad f\in \mathcal{M}_{\mathcal{Q}_{\mathcal{Z}}}$$ is contractive.
     \end{theorem}

\begin{proof}
% Since the points $\boldsymbol\zeta_1,\ldots,\boldsymbol\zeta_p$ are distinct, the corresponding kernel functions are linearly independent, and hence the Gram matrix $G_{\mathcal Z}$ is invertible.
    Consider the operator $X_{\mathcal{Z,W}}\in \mathcal{B}(\mathcal{Q}_{\mathcal{Z}})$ defined by \beq \label{eq 12} X_{\mathcal{Z,W}}^* \mathbb S_{\boldsymbol \theta(\Omega)}(\cdot, \boldsymbol \zeta_i)=\overline{w}_i \mathbb S_{\boldsymbol \theta(\Omega)}(\cdot, \boldsymbol \zeta_i),\qquad i=1,\ldots,p.\eeq 
    A routine verification shows that $X_{\mathcal{Z},\mathcal{W}}$ is a bounded module map. Hence, by our preceding discussion, $\psi=X_{\mathcal{Z,W}}P_{\mathcal{Q}_{\mathcal{Z}}}1\in H^\infty(\boldsymbol \theta(\Omega))$ and $X_{\mathcal{Z,W}}=S_{\psi}$ on $\mathcal{Q}_{\mathcal{Z}}.$

    Suppose there exists $\varphi\in \mathcal{S}(\boldsymbol \theta(\Omega))$ such that $\varphi(\boldsymbol{\zeta}_i) = w_i$ for all $i = 1, \ldots, p$. Using the reproducing property,
    $$S_{\varphi}^* \mathbb S_{\boldsymbol \theta(\Omega)}(\cdot, \boldsymbol \zeta_i)=P_{\mathcal{Q}_{\mathcal{Z}}} T_{\varphi}^*|_{\mathcal{Q}_{\mathcal{Z}}} S_{\boldsymbol \theta(\Omega)}(\cdot, \boldsymbol \zeta_i) = \overline{\varphi(\boldsymbol \zeta_i)}\mathbb S_{\boldsymbol \theta(\Omega)}(\cdot, \boldsymbol \zeta_i)=\overline{w}_i \mathbb S_{\boldsymbol \theta(\Omega)}(\cdot, \boldsymbol \zeta_i).$$ 
    Comparing this with \eqref{eq 12}, we obtain $X_{\mathcal{Z,W}}=S_{\varphi}.$ Thus, $X_{\mathcal{Z},\mathcal{W}}$ admits a lift to the Schur function $\varphi.$
    
    Conversely, if $X_{\mathcal{Z,W}}=S_{\varphi}$ for some $\varphi\in \mathcal{S}(\boldsymbol \theta(\Omega)),$ then \eqref{eq 12} gives
    $$\overline{w}_i \mathbb S_{\boldsymbol \theta(\Omega)}(\cdot, \boldsymbol \zeta_i)=\overline{\varphi(\boldsymbol \zeta_i)} \mathbb S_{\boldsymbol \theta(\Omega)}(\cdot,\boldsymbol \zeta_i), \qquad i=1,\ldots,p,$$ and hence  $\varphi(\boldsymbol \zeta_i)=w_i$ for all $i=1,\ldots,p.$

    Therefore, the interpolation problem is solvable if and only if the operator $X_{\mathcal{Z},\mathcal{W}}$ admits a lift to the Schur-class interpolant.
   % Therefore the interpolation problem is equivalent to the liftability of $X_{\mathcal{Z},\mathcal{W}}$ to the interpolant. 
   % there exists $\varphi\in \mathcal{S}(\boldsymbol \theta(\Omega))$ satisfying  $\varphi(\boldsymbol \zeta_i)=w_i$ for all $i=1,\ldots,p$ if and only if the operator $X_{\mathcal{Z},\mathcal{W}}$ admits a lift to the Schur function $\varphi$.
    By Theorem \ref{1st characterization}, this occurs if and only if the functional 
    $X_{\mathcal{Q}_{\mathcal{Z}}}: \left(\mathcal{M}_{\mathcal{Q}_{\mathcal{Z}}}, \|\cdot \|_{L^1(\partial\boldsymbol \theta(\Omega))}\right) \to \mathbb C$ defined by $$X_{\mathcal{Q}_{\mathcal{Z}}}(f)=\int_{\partial\boldsymbol \theta(\Omega)} \psi f d\mu_{\boldsymbol \theta}, \qquad f\in \mathcal{M}_{\mathcal{Q}_{\mathcal{Z}}},$$
    is contractive.
    Moreover, in this contractive setting, $X_{\mathcal{Z}, \mathcal{W}} = S_{\psi} = S_{\varphi}$ on $\mathcal{Q}_{\mathcal{Z}}$, and $\psi(\boldsymbol{\zeta}_i) = w_i$ for all $i = 1, \ldots, p.$
    
    It remains only to identify $\psi.$ Since $\psi\in \mathcal{Q}_{\mathcal{Z}},$ there are scalars $c_1,\ldots,c_p$ such that $\psi=\sum_{j=1}^p c_j \mathbb S_{\boldsymbol \theta(\Omega)}(\cdot, \boldsymbol \zeta_j).$ Evaluating at each node $\boldsymbol \zeta_i$ yields
    $$ w_i=\psi(\boldsymbol \zeta_i)= \sum_{j=1}^p c_j \mathbb S_{\boldsymbol \theta(\Omega)}(\boldsymbol \zeta_i, \boldsymbol \zeta_j),\qquad i=1,\ldots,p.$$
    This uniquely determines the coefficients $c_1,\ldots, c_p$, and hence  $\psi=\psi_{\mathcal{Z,W}}.$
    \end{proof}

\end{document}
